\chapter{Resumen}\label{cha:resumen}

	Las radiogalaxias son un tipo de galaxia activa que se caracterizan por presentar intensa emisi\'on en radiofrecuencias. La misma se origina en \textit{jets} relativistas impulsados por un agujero negro supermasivo central. 
	Observaciones recientes de muy alta resoluci\'on, como las obtenidas mediante interferometr\'ia de muy larga base (VLBI, por sus siglas en ingl\'es), han revelado estructuras complejas en estos \texit{jets}, incluyendo variaciones espaciales significativas en sus propiedades f\'isicas.

	En este trabajo, desarrollamos un modelo teórico de \textit{jet} relativista multizona con estructura lateral, con el objetivo de interpretar la morfología observada en fuentes como M87 y Centaurus A. 
	Para ello, implementamos c\'alculos num\'ericos del transporte de part\'iculas a lo largo del \textit{jet} y de la radiaci\'on sincrotr\'on de las mismas. 
	Calculamos distribuciones espectrales de energía y mapas sintéticos de intensidad para la fuente de Centaurus A y los contrastamos con las más recientes observaciones del Event Horizon Telescope. 
	Esto nos permite caracterizar las propiedades espaciales del \textit{jet} y comprender mejor la f\'isica detr\'as del lanzamiento de los mismos.





\newpage

\chapter{Abstract}\label{cha:abstract}

	Radio galaxies are a type of active galaxy characterized by intense emission at radio frequencies. This emission originates in relativistic \textit{jets} driven by a supermassive black hole at the center. 
	Recent high-resolution observations, such as those obtained using very long baseline interferometry (VLBI), have revealed complex structures in these \textit{jets}, including significant spatial variations in their physical properties. 
	
	In this work, we develop a theoretical model of a multi-zone relativistic \textit{jet} with lateral structure, with the aim of interpreting the observed morphology in sources such as M87 and Centaurus A. 
	To this end, we implement numerical calculations of particle transport along the \textit{jet} and the synchrotron radiation from {said particles}. We calculate spectral energy distributions and synthetic intensity maps for the Centaurus A source and contrast {them} with the most recent observations from the Event Horizon Telescope. 
	This allows us to characterize the spatial properties of the \textit{jet} and better understand the physics behind their launch.

\endinput
